\documentclass[a4paper,11pt]{article}

\usepackage[ngerman]{babel}
\usepackage[top=3cm,bottom=3cm,left=3cm,right=3cm]{geometry} %NICHT IN KEMIE2


%\usepackage{microtype} %schöneres typesetting  mit [stretch=10] für nur 1% anstelle von 2%
%\usepackage{lmodern} %Latin Modern FONT
%\usepackage{pifont} %schönere Font



\usepackage{amsmath}
\usepackage{nicefrac}
\usepackage{amssymb}
%\usepackage{mathtools}

\usepackage{titlesec} %Um Section, etc bearbeiten zu können
\usepackage{setspace}
\usepackage{enumitem} %enumerate
\usepackage{multicol}
\usepackage[many]{tcolorbox}
\usepackage{hyperref}
\usepackage{arydshln}

%\renewcommand{\ttdefault}{lmtt} %Schriftart wird geändert zu Latin Modern Typewriter


\hypersetup{hidelinks,
			colorlinks=true,
			linkcolor=black,
			citecolor=miudiutex,
			urlcolor=miugray2
			}

\setlength{\parindent}{0pt}
\setcounter{section}{0}


\titlespacing*{\section}{0pt}{20mm}{6mm}
\titlespacing*{\subsection}{0pt}{6mm}{3mm}
\titlespacing*{\subsubsection}{0pt}{3mm}{0mm}



%\definecolor{miudiutex}{HTML}{2f3d39}
\definecolor{miudiutex}{HTML}{1d2926}
\definecolor{miugray}{HTML}{b4cccf}
\definecolor{miugray2}{HTML}{4d7365}
\definecolor{red}{HTML}{cf1f5c}
\definecolor{green}{HTML}{00A36C}
\definecolor{delphinidin}{HTML}{315794}
\definecolor{pelargonidin}{HTML}{b33807}
\definecolor{cyanidin}{HTML}{ba0746}
\definecolor{petunidin}{HTML}{b56fed}



\raggedcolumns %wichtig für Multicolumns


%================================
% SYMBOLS
%================================

\newcommand{\RA}{$\rightarrow$~}
\newcommand{\RL}{$\rightleftharpoons$~}
\newcommand{\HARP}{\ding{226}~}
%$\xrightarrow{2}$ %Pfeil mit 2 drüber


%================================
% SHORT COMMANDS (BOXES ; ETC)
%================================


\newcommand{\notiz}[1]{\begin{notizz}{}{}\begin{center}
			{\color{miugray2!50!miudiutex}#1} \end{center}\end{notizz}}
\newcommand{\zit}[2]{\begin{zitat*}{#1}{}\small #2 \end{zitat*}}

\newcommand{\frage}[2]{\begin{qst*}{#1}{}\small #2 \end{qst*}}

\newcommand{\unten}{\end{center} \tcblower \begin{center}}


\newcommand*\circled[1]{\tikz[baseline=(char.base)]{
            \node[shape=circle,draw,inner sep=0.5pt,miudiutex] (char) {\tiny #1};}}
            
\newcommand{\miusquare}{{\color{miugray2!50}\scriptsize $\blacksquare$}~}
\newcommand{\miusquarp}{{\color{petunidin!50}\scriptsize $\blacktriangleright$}~}
\newcommand{\niu}{\item[\miusquare]}
\newcommand{\nip}{\item[\miusquarp]}

%%%%% ABSAETZE

\newcommand{\pps}{\par \vspace*{1mm}}
\newcommand{\pp}{\par \vspace*{3mm}}
\newcommand{\ppbig}{\par \vspace*{8mm}}
\newcommand{\pphuge}{\par \vspace*{10mm}}

%================================
% FRAGE BOX
%================================

\tcbuselibrary{theorems,skins,hooks}
\newtcbtheorem{qst}{Frage:}
{%
	enhanced,
	breakable,
	colback = gray!15!white,
	frame hidden,
	boxrule = 0sp,
	borderline west = {2.5pt}{0pt}{red!70},
	sharp corners,
	detach title,
	before upper = \tcbtitle\par\smallskip,
	coltitle = black,
	fonttitle = \bfseries\sffamily,
	description font = \mdseries,
	separator sign none,
	segmentation style={solid, gray!50!black},
}
{th}




%================================
% ZITAT BOX
%================================

\tcbuselibrary{theorems,skins,hooks}
\newtcbtheorem{zitat}{Zitat}
{%
	enhanced,
	breakable,
	colback = gray!15!white,
	frame hidden,
	boxrule = 0sp,
	borderline west = {2.5pt}{0pt}{black!70},
	sharp corners,
	detach title,
%	before upper = \tcbtitle\par\smallskip,
	coltitle = gray!50!black,
	fonttitle = \bfseries\sffamily,
	description font = \mdseries,
%	separator sign none,
	segmentation style={solid, gray!50!black},
}
{th}



%================================
% SIMPLE SCHWARZER RAHMEN BOX
%================================


\newcommand{\boxfr}[1]{
	\begin{tcolorbox}[
	drop shadow southeast,
	enhanced,
	colframe=black!50,
	colback=white,
	boxrule = 1pt, 
%	toprule=0pt, bottomrule=0pt,rightrule=0pt,leftrule=0pt,
	boxsep=5pt,
	arc=1pt,
	]
	\begin{center}
	#1
	\end{center}
	\end{tcolorbox}
}

\definecolor{miudiutex}{HTML}{1d2926}
\definecolor{miugray}{HTML}{b4cccf}
\definecolor{miugray2}{HTML}{4d7365}
\definecolor{white2}{HTML}{f5f7f8}
\definecolor{red}{HTML}{cf1f5c}
\definecolor{green}{HTML}{00A36C}
\definecolor{delphinidin}{HTML}{315794}
\definecolor{uniblau}{HTML}{004e9f}
\definecolor{unigelb}{HTML}{fcba00}
\definecolor{unigrau}{HTML}{909085}
\definecolor{pelargonidin}{HTML}{b33807}
\definecolor{cyanidin}{HTML}{ba0746}
\definecolor{paeonidin}{HTML}{d15ca6}
\definecolor{petunidin}{HTML}{b56fed}
\definecolor{malvidin}{HTML}{8a2d8c}

\newcommand{\metermilli}{\mskip3mu \text{mm}}
\newcommand{\cm}{\mskip3mu \text{cm}}
\newcommand{\m}{ \text{m}}
\newcommand{\lite}{ \text{l}}
\newcommand{\kg}{ \text{kg}}
\newcommand{\g}{ \text{g}}
\newcommand{\km}{ \text{km}}
\newcommand{\h}{ \text{h}}
\newcommand{\s}{ \text{s}}
\newcommand{\tonne}{ \text{t}}
\usepackage[utf8]{inputenc}
\usepackage[T1]{fontenc}
\usepackage{lmodern}

\usepackage[
    activate={true,nocompatibility},
    final,
    tracking=true,
    kerning=true,
    spacing=true,
    factor=1100,
    stretch=30,
    shrink=20
]{microtype}
\usepackage{pdfpages}
\usepackage{awesomebox}
\usepackage{marginnote}
\tcbuselibrary{documentation}
\usepackage{sidenotes}
\usepackage{import}
\usepackage{xifthen}
\usepackage{transparent}
\usepackage[table]{xcolor}


\newcommand{\abbicon}{%
  \protect\tikz[baseline=0.2mm,scale=0.8]{%
    % Das blaue Quadrat mit abgerundeten Ecken
    \fill[uniblau, rounded corners=1.2pt] (0,0) rectangle (0.8em, 0.8em);
    % Der weiße Kreis in der Mitte
    \fill[white] (0.4em, 0.4em) circle (0.12em);
  }%
} 
\usepackage[labelfont={bf},font={color=miudiutex,small},figurename={\abbicon $~$Abbildung}]{caption}


\newcommand{\bckgrnd}{
\AddToHookNext{shipout/background}{
\begin{tikzpicture}[remember picture, overlay]
 \draw[draw=none,fill=uniblau!70, shift={(current page.north west)}] (-3,0) rectangle (23,-3.75);
 \draw[draw=none,fill=miugray, shift={(current page.north east)}] (-7,0) rectangle (3,-3.75);
\end{tikzpicture}
}
}



\newcommand{\incfig}[2]{%
\def\svgwidth{#2\columnwidth}
\import{C://LaTeX-Files/pngs/}{#1.pdf_tex}
}

\renewcommand{\textbf}[1]{\textsf{{\bfseries {#1}}}}



\titleformat{\section}[hang]
		{\LARGE \color{uniblau!75!black} \sffamily \bfseries}
		{\thesection}
		{6mm}
		{}
		[]
\titleformat{\subsection}[hang]
		{\large \sffamily \bfseries \color{black}}
		{\thesubsection}
		{4mm}
		{{\color{miugray}$\blacksquare ~$}}

\titleformat{\subsubsection}[block]
		{\normalsize \bfseries \sffamily \color{miudiutex}}%
		{\normalsize \color{black} \thesubsubsection}
		{2mm}
		{{\color{miugray}$\blacksquare ~$}}
		[ \color{black}]
		
\titleformat{\paragraph}
		{\normalsize \bfseries \sffamily \color{black}}
		{\footnotesize \theparagraph}
		{1mm}
		{}

\pagestyle{empty}
\titlespacing*{\section}{0pt}{20mm}{12mm}
\titlespacing{\section}{0pt}{20mm}{12mm}

\titlespacing*{\subsection}{0pt}{14mm}{4mm}
\titlespacing{\subsection}{0pt}{14mm}{4mm}

\titlespacing*{\subsubsection}{0pt}{6mm}{3mm}
\titlespacing{\subsubsection}{0pt}{6mm}{3mm}

\titlespacing*{\paragraph}{0pt}{6mm}{2mm}
\titlespacing{\paragraph}{0pt}{6mm}{2mm}


\let\oldmarginfigure\marginfigure
\let\endoldmarginfigure\endmarginfigure

\let\oldmarginfigure\marginfigure
\let\endoldmarginfigure\endmarginfigure

\renewenvironment{marginfigure}[1][0pt] % [1] steht für 1 Argument, [0pt] ist der Standardwert
  {\oldmarginfigure[#1]\captionsetup{labelfont={sf,bf,color=uniblau!75!black},font={color=miudiutex,small}}}
  {\endoldmarginfigure}
  
  
  
  
\renewcommand{\labelitemi}{\color{uniblau!70}$\bullet$}
\renewcommand{\labelitemii}{\color{uniblau!70}$\cdot$}
\newcommand{\plmtsquare}{{\color{uniblau!70}$\blacksquare~$}}
\newcommand{\splmt}[1]{{\color{uniblau!75!black} \sffamily #1}}
\newcommand{\fplmt}[1]{{\color{uniblau!75!black} \sffamily \bfseries #1}}

\newcommand{\antwort}[1]{
\begin{tcolorbox}[
	enhanced,
	enforce breakable,
	standard jigsaw,
	boxrule=0.7pt,
	colframe=black,
	sharp corners,
	boxsep=5pt,
	title=\bfseries \sffamily Antwort,
	opacityback=0,
	coltitle=miudiutex,
	attach boxed title to top text left={yshift=-\tcboxedtitleheight/2},
	boxed title style={colback=white,colframe=white},
	bottomrule at break=0pt,
	toprule at break=0pt,
	pad at break=16mm,
]
	\begin{center}
	#1
	\end{center}
	\end{tcolorbox}
}




\rowcolors{2}{miugray!50}{miugray!50}
\arrayrulecolor{white} % Weiße Gitterlinien
\setlength{\arrayrulewidth}{1.5pt} % Etwas dickere Linien
\renewcommand{\arraystretch}{1.4} % Mehr Zeilenhöhe für bessere Lesbarkeit}


%\setlist[itemize]{itemsep=0mm}
\setlist[enumerate]{itemsep=0mm}


\newif\ifshowme
%\showmetrue




\begin{document}
\newgeometry{top=1.8cm,bottom=4cm,left=1.5cm,right=7.5cm,marginparsep=-2.00cm,marginparwidth=4.75cm}


\bckgrnd


\begin{center}
\LARGE
\color{white}
\textbf{Übung 04} \par \vspace*{2mm}
\large \color{miugray} \textbf{Prozessbezogene Lebensmitteltechnologie}
\end{center}
\par 
\vspace*{6mm}


\color{miudiutex}
\section*{Wärmeübertragung I}
\subsection*{Aufgabe 1}
Berechnen Sie die Strahlungsleistung, die von 100 m$^2$ einer polierten Eisenoberfläche
(Emissionsgrad = 0,06) abgegeben wird. Die Temperatur der Oberfläche beträgt 37 °C.

%\begin{marginfigure}
%\centering
%\includegraphics[width=4.5cm]{F://LaTeX-Files/pngs/strawberry}
%\caption{Erdbeerbecher mit Vanilleschmand und Gewürzcrumble}
%\end{marginfigure}

\ifshowme
\antwort{
\begin{align*}
P &= \varepsilon \cdot \sigma \cdot A \cdot T^4 \\ 
P &= 0,06 \cdot  100~\text{m}^2 \cdot 5,67 \cdot 10^{-8} ~\text{W} \cdot \text{m}^{-2}\text{K}^{-4}\cdot 310,15~\text{K}^4 \\
P &= 3.147,9~\text{W}
\end{align*}
}
\else
\fi

\subsection*{Aufgabe 2}
Eine Seite einer 1 cm dicken Edelstahlplatte wird auf 110 °C gehalten, die andere Seite
auf 90 °C. Berechnen Sie unter der Annahme stationärer Bedingungen die
Wärmeübertragungsrate pro Flächeneinheit durch die Platte. Die Wärmeleitfähigkeit von
nichtrostendem Stahl beträgt \mbox{17 W m$^{-1}$ °C$^{-1}$}.

\ifshowme
\antwort{
\begin{align*}
\dot{Q} &= - A \cdot \frac{\lambda}{d} \cdot (T_2 - T_1) \\
- \frac{\dot{Q}}{A} &= \cdot \frac{\lambda}{d} \cdot (T_2 - T_1) \\
- \frac{\dot{Q}}{A} &= \cdot \frac{17 \text{W} / \text{m°C}}{0,01m} \cdot 20\text{°C}\\
	&= 34.000 ~\text{W / m$^2$}
\end{align*}

}
\pagebreak
\else
\fi

\subsection*{Aufgabe 3}
Ein Glühofen hat eine 4 m$^2$ große Wand aus \splmt{Schamotte}. Die Dicke der Wand sei 0,5 m.
Zwischen der Innen- und der Außenseite der Wand besteht eine Temperaturdifferenz von
1.000 K. Die Temperatur außerhalb des Ofens beträgt 20 °C. Die Wärmeleitfähigkeit der
Schamotte ist $\lambda_1$ = 0,93 W/mK. 


\marginnote[]{\small Als \fplmt{Schamotte} werden im allgemeinen Sprachgebrauch häufig alle feuerfesten Steine und Ausmauerungen bezeichnet.}[-2.5cm]

\begin{enumerate}[label=\alph*)]
\item Wie groß ist der tägliche Wärmeverlust durch die Wand?


\ifshowme
\antwort{
\begin{align*}
\dot{Q} &= - A \cdot \frac{\lambda}{d} \cdot (T_2 - T_1) \\
&= - 4m^2 \cdot \frac{0,93}{0,5m} \cdot 1000 K \\
&= - 7.440~W \\
&= - 7.440~W \cdot 24h \\
&= - 178.560~Wh
\end{align*}
}
\else
\fi

\item Um wie viel Prozent verringert sich dieser Verlust, wenn die Außenwand (von außen)
durch eine 3 cm dicke Isolierschicht mit einer Wärmeleitfähigkeit von $\lambda_2$ = 0,093 
W/mK verstärkt wird?


\ifshowme
\antwort{
\begin{align*}
\dot{Q} &= - \frac{A \cdot \Delta T_{ges}}{(\frac{d_2}{\lambda_2}+\frac{d_1}{\lambda_1})} \\
\end{align*}
Einsetzen:
\begin{align*}
\dot{Q} &= - \frac{4m^2 \cdot 1000 K}{(\frac{0,03m}{0,093 W/mk}+\frac{0,5m}{0,93 W/mk})} \\
&= 4.650 W
\end{align*}
Prozentualer Unterschied:
\begin{align*}
\frac{4.650 W}{7.440 W} &\approx  0,625 = 62,5\%
\end{align*}
Also verringert sich der Verlust um 100-62,5\% = 37,5\%

}
\pagebreak
\else
\fi

\item Welche Temperatur herrscht an der Grenze zwischen der Isolierung und der
Schamotte-Schicht?

\ifshowme
\begin{marginfigure}[0cm]
\else
\begin{marginfigure}[-4cm]
\fi

\centering
\begin{tikzpicture}
\draw[fill=miugray!50] 	(0,0) rectangle (0.8,3);
\draw[fill=white]		(0.8,0) rectangle (2,3);
\node at (0.4,1.5) {$\dot{Q_1}$};
\node at (1.4,1.5) {$\dot{Q_2}$};
\draw[->,shorten >=2pt] 	(0,-0.3) -- (0,0)	node[below,yshift=-2mm]{$T_1$}; 
\draw[->,shorten >=2pt] 	(0.8,-0.3) -- (0.8,0)node[below,yshift=-2mm]{$T_2$};
\draw[->,shorten >=2pt] 	(2,-0.3) -- (2,0)	node[below,yshift=-2mm]{$T_3$};
\end{tikzpicture}
\caption{Der Wärmestrom ist an jeder Stelle der Wand konstant ($\dot{Q_1} = \dot{Q_2}$). Die Differenz der Temperatur $\Delta T$ lässt sich in ihre Bestandteile $T_1$, $T_2$ und $T_3$ zerlegen.}
\end{marginfigure}

\ifshowme
\antwort{
\begin{align*}
\dot{Q} &= - \lambda \cdot A \cdot \frac{T_3 - T_2}{d} \\
T_2 &= \frac{\dot{Q} \cdot d}{\lambda \cdot A } + T_3 \\
\\
T_2 &= \frac{4.650 ~\text{W} \cdot 0,03~\text{m}}{0,093 ~\text{W/mk} \cdot 4~\text{m}^2} + 20~^{\circ}\text{C} \\
 &= 395~^{\circ}\text{C}
\end{align*}
}

\else
\fi

\end{enumerate}


\subsection*{Aufgabe 4}
\marginnote[]{\small Für die Berechnung der spezifischen \fplmt{Wärmekapazität $c$} kann folgende Formel verwendet
werden:

\begin{align*}
c = &\quad ~   4,180 \times \text{Wasser} \\& + 1,711 \times \text{Proteine} \\& + 1,928 \times \text{Fett} \\&+ 1,547 \times \text{Kohlenhydrate} \\& + 0,908 \times \text{weitere Bestandteile}
\end{align*}
}

Welche Wärmemenge (in kWh) benötigen Sie, um 6,5 L einer Kartoffelsuppe von 10 °C auf
60 °C zu erhitzen? Es wird angenommen, dass die Dichte der Suppe 1 kg/dm$^3$ beträgt und
es keine Wärmeverluste gibt. \pp 

Die Suppe setzt sich pro 100 g zusammen aus: \par 
\begin{itemize}
\item Fett: 14,3 g
\item Proteine: 2,7 g
\item Kohlenhydrate: 11,5 g
\item Weitere Bestandteile: 0,5 g
\end{itemize}




\ifshowme
\antwort{
Zunächst Wärmekapazität $c$ ausrechnen:
\begin{align*}
c &= ~~~ 4,180 \cdot 0,71 + 1,711 \cdot 0,027 + 1,928 \cdot 0,143 \\
&~~~ + 1,547 \cdot 0,115 + 0,908 \cdot 0,005 \\
 &= ~~~ 3,472146~ \text{kJ} \cdot \text{kg}^{-1} \cdot \text{K}^{-1}
\end{align*}
Energie berechnen:
\begin{align*}
Q &= m \cdot c \cdot \Delta T \\ \\
Q &= 6,5~ \text{kg} \cdot 3,472~ \text{kJ} \cdot \text{kg}^{-1}\cdot \text{K}^{-1} \cdot (333,15-283,15 ~ \text{K}) \\
&= 1.128,45~ \text{kJ} \\
&= 1.128,45~ \text{kWs}  \\
&= 0,313~ \text{kWh}
\end{align*}

}
\pagebreak
\else
\fi



\subsection*{Aufgabe 5}
In Ihrem 20 °C warmen Zuhause schalten Sie Ihre Herdplatte auf maximaler Stufe an. Diese erhitzt sich daraufhin auf 300 °C. Im nächsten Schritt stellen Sie eine Pfanne aus Aluminium ($\lambda_{Al} = 236~\text{W/mK}$) mit einem Durchmesser von 28 cm und einer Dicke von 2 mm auf die Herdplatte.
\begin{itemize}
\item Welcher Wärmefluss entsteht zwischen der Unter- und der Oberseite der Pfanne?
\ifshowme
\antwort{
\begin{align*}
\dot{Q} &= \lambda_{Al} \cdot \frac{A}{d} \cdot (\Delta T) \\
&= 236~\text{W/mK} \cdot \frac{\pi \cdot 0,14~\text{m}^2}{0,002~\text{m}} \cdot 280~\text{K}\\
&= 2.034.445~\text{J}\\
&= 2.034~\text{kJ}
\end{align*}
}
\else
\fi

\item Wieso können Sie von dem fließenden Wärmestrom keine Rückschlüsse auf ihre tatsächliche Stromrechnung schließen?
\ifshowme
\antwort{
Das Delta T wird mit der Zeit immer geringer. Die Herdplatte muss anfänglich eine viel größere Leistung aufbringen um die kalte Herdplatte (bzw auch die Pfanne) zu erhitzen. Ein typischer Herd hat eine Leistung von circa 10 kW. Für einen konstanten Wärmestrom aus unserer Aufgabe wären 2 Megawatt notwendig!
}
\pagebreak
\else
\fi

\end{itemize}



\subsection*{Aufgabe 6}
Eine gefrorene Kartoffelsuppe (-19 °C, 850 g) soll mithilfe einer Mikrowelle zum Verzehr auf
80 °C erwärmt werden. Die Leistung der Mikrowelle liegt bei 800 W. Nehmen Sie an, dass
diese Leistung zu 100\% für die Erwärmung der Kartoffelsuppe zur Verfügung steht.
Die spezifische Wärmekapazität der gefrorenen Suppe beträgt 2,0 kJ/(kgK) und die der
aufgetauten Suppe beträgt 4,62 kJ/(kgK). Es werden eine spezifische Schmelzwärme von 334
kJ/kg und eine Dichte der Suppe von 1 g/cm$^3$ angenommen. \pp
Wie lange dauert es, bis die gewünschte Temperatur erreicht ist?

\ifshowme
\antwort{
\begin{align*}
Q &= m \cdot c \cdot \Delta T \\
Q &= 0,85~\text{kg} \cdot 2 ~\text{kJ/kgK} \cdot 19 ~\text{K} + 0,85~\text{kg} \cdot 334 ~\text{kJ/kg} 
\\&~~~ + 0,85~\text{kg} \cdot 4,62 ~\text{kJ/kgK} \cdot 80 ~\text{K}\\
	&= 630,36 ~\text{kJ}  \\
	&= 630.360,0 ~\text{Ws } && | : 800 ~\text{W} \\
	&= 787.95~\text{s} \\
	&\approx 13~\text{min}~ 8~\text{s}
\end{align*}

}
\else
\pagebreak
\fi






\subsection*{Aufgabe 7}
Auf welche Temperatur kühlt sich ein Getränk (500 g, 20 °C) ab, wenn 50 g Eis mit einer
Temperatur von -20 °C hinzugegeben werden? Die spezifische Wärmekapazität von Eis liegt
bei 2,0 kJ/(kgK) und von Wasser bei 4,2 kJ/(kg K). Die spezifische Schmelzwärme von Eis ist angegeben mit $q_S$ = 334 kJ/kg. Vernachlässigen Sie etwaige Wärmeströme in Bezug auf das
Glas oder die Luft. \pp
Berechnen Sie die Endtemperatur $T_f$. Überlegen Sie dafür zunächst, in welche (berechenbare) Schritte Sie diesen Prozess aufteilen können.


\ifshowme
\antwort{
Zunächst Energie des Eiswürfels ausrechnen (wir gehen davon aus, dass er schmelzen wird)


\begin{align*}
Q &= m \cdot c \cdot \Delta T \\ 
Q &= 0,05 ~\text{kg} \cdot 2 ~\text{kJ/kg} \cdot 20 ~\text{K} + 0,05 \cdot 334 ~\text{kJ/kg} \\ 
Q &= 18,7 ~\text{kJ} 
\end{align*}

Im nächsten Schritt schauen wir uns an, wie sehr wie die Temperatur unseres Getränkes mit der Energie aus dem Eiswürfel absenken können:

\begin{align*}
Q &= m \cdot c \cdot \Delta T \qquad \rightarrow \qquad \Delta T = \frac{Q}{m \cdot c} \\ \\
\Delta T &= \frac{18,7 ~\text{kJ}}{0,5 ~\text{kg} \cdot 4,2 ~\text{kJ/kgK}} \\ 
		&= 8,9 ~\text{°C}
\end{align*}

Wir wissen nun, dass unser Getränk eine Temperatur von 11,1°C (20-8,9 °C) und unsere Flüssigkeit, die ehemals Eiswürfel war, eine Temperatur von 0°C hat. Wenn zwei Flüssigkeiten mit unterschiedlichen Temperaturen gemischt werden, gilt für die Endtemperatur: 


\begin{align*}
T_f &= \frac{c_1 \cdot m_1 \cdot T_1 + c_2 \cdot m_2 \cdot T_2}{c_1 \cdot m_1 + c_2 \cdot m_2} \\
T_f &= \frac{m_1 \cdot T_1 + m_2 \cdot T_2}{m_1 + m_2} \\
T_f &= \frac{0,5~\text{kg} \cdot 11,1 ~\text{°C} + 0,05~\text{kg} \cdot 0 ~\text{°C}}{0,5 ~\text{kg} + 0,05 ~\text{kg}} \\
	&= 10,09 ~\text{°C}
\end{align*}

}

\else
\fi



%\subsection*{Aufgabe 8}
%Für eine Mikrowellenfrequenz von 2.450 MHz und eine Temperatur von \mbox{20 °C} beträgt die
%stoffabhängige Permitivitätszahl $\varepsilon '$ von rohen Kartoffeln 64 und der Verlustwinkel tan $\delta$ 0,23.
%Nehmen Sie als Lichtgeschwindigkeit den gerundeten Wert 3 $\cdot$ 10$^8$ m/s an.
%\begin{enumerate}[label=\alph*)]
%
%\item Bestimmen Sie anhand untenstehender Formel die Eindringtiefe des Feldes in die
%Kartoffel.
%\item Bestimmen Sie die Eindringtiefe auch für Wasser bei 25°C ($\varepsilon '$ 78; tan $\delta$: 0,16) und Eis
%($\varepsilon '$: 3,2; tan $\delta$: 0,0009).
%\item Wie ist die Eindringtiefe definiert?
%
%\end{enumerate}
%Es gilt für die Eindringtiefe Z mit der Wellenlänge $\lambda$: 
%\begin{align*}
%Z = \frac{\lambda}{2 \pi} \left[ \frac{2}{\varepsilon ' \cdot (\sqrt{1+ \tan^2 \delta - 1)}} \right]^{\frac{1}{2}}
%\end{align*}
%
%\ppbig 
%
%
%\ifshowme
%\antwort{
%Für das Lösen der oben stehenden Aufgaben müssen wir zunächst die Wellenlänge $\lambda$ ausrechnen:
%
%
%\begin{align*}
%\lambda = \frac{c}{f} \qquad \rightarrow \qquad \frac{3 \cdot 10^8 ~\text{m/s}}{2.450 \cdot 10^6 ~\text{Hz}} = 0,1224 ~\text{m}
%\end{align*}
%
%\marginnote[]{Vorsicht beim Einsetzen von tan $\delta$. Die Formel ist einwenig verwirrend.}[1cm]
%Dann Werte einsetzen (Kartoffel): 
%
%\begin{align*}
%Z &= \frac{0,1224 ~\text{m}}{2 \pi} \left[ \frac{2}{64 \cdot (\sqrt{1+ (0,23)^2 - 1)}} \right]^{\frac{1}{2}} \\
% &= 0,0213 ~\text{m}
%\end{align*}
%
%Das Selbe für Wasser: 
%
%\begin{align*}
%Z &= \frac{0,1224 ~\text{m}}{2 \pi} \left[ \frac{2}{78 \cdot (\sqrt{1+ (0,16)^2 - 1)}} \right]^{\frac{1}{2}} \\
% &= 0,0276 ~\text{m}
%\end{align*}
%
%
%Und Eis:
%
%\begin{align*}
%Z &= \frac{0,1224 ~\text{m}}{2 \pi} \left[ \frac{2}{2,3 \cdot (\sqrt{1+ (0,0009)^2 - 1)}} \right]^{\frac{1}{2}} \\
% &= 24,2 ~\text{m}
%\end{align*}
%
%
%}
%\else
%\fi
%
%
%
%
%\ifshowme
%\antwort{
%\textbf{Eindringtiefe Z:} \par Punkt, an dem die elektrische Feldstärke nur noch 1/e von der ursprünglichen Feldstärke entspricht.
%
%}
%\else
%\fi




\end{document}
